\begin{justify}
Research on effort and time estimation shows that estimation accuracy is an important factor in software and project management. Boehm's work on software engineering economics established the importance of estimating effort using measurable project attributes and historical knowledge \cite{boehm1981software}. Although the present project is smaller than industrial cost estimation models, it follows the same basic idea that better records can improve future planning.

Jorgensen and Shepperd reviewed several software development cost estimation studies and observed that estimation methods differ in accuracy depending on available data, project context, and estimator experience \cite{jorgensen2007forecasting}. This supports the need for systems that collect actual performance data instead of relying only on memory or guesswork.

Molokken and Jorgensen reviewed surveys on software effort estimation and highlighted common estimation problems such as optimism, uncertainty, and weak historical tracking \cite{molokken2003review}. The proposed system addresses these issues at a personal task level by storing expected and actual time for each task and providing direct feedback.

Machine learning methods such as recurrent neural networks and Long Short-Term Memory models can be used for sequence-based prediction when a large amount of historical data is available \cite{hochreiter1997lstm}. However, the current project uses rule-based artificial intelligence because it is transparent, simple to validate, and suitable for a mini project. In future, collected task history can be used to train a prediction model for suggested completion time.
\end{justify}
